\section{Conclusão}
\label{sec:conclusao}

Este trabalho caracterizou longitudinalmente a evolução da qualidade estrutural de
projetos Java do ecossistema Spring criados até 2017 e ativos na era contemporânea
(E4, \texttt{pushed\_at >= 2022-01-01}), acompanhando versões históricas com
rastreabilidade estrita por commit Git até outubro de 2026. Em substituição a filtros
arbitrários de popularidade por número de estrelas, a relevância de engenharia dos 785
candidatos iniciais foi aferida localmente pela aplicação simultânea dos quatro cenários
do RepoReaper (C1 de finalidade não didática, C2 ajustado para projetos brinquedo com
flexibilização dos termos \texttt{demo}, \texttt{sample} e \texttt{example} próprios de
convenções do Spring Boot, C3 de automação de build Maven/Gradle e C4 de presença de
testes automatizados), aprovando 607 repositórios. A exigência metodológica de dados CK
válidos em todas as eras previstas para cada repositório consolidou a amostra final de
\textbf{470 repositórios válidos (1.427 snapshots CK)}, organizados em 17 projetos do
grupo pioneiro (2009--2013, acompanhados continuamente nas quatro eras E1--E4) e 453
projetos do grupo principal (2014--2017, acompanhados continuamente nas três eras
E2--E4).

\textbf{Objetivo (i), trajetórias estruturais.} O acoplamento entre objetos (CBO) e o
conjunto de resposta das classes (RFC) cresceram de forma monotônica nos mesmos projetos
ao longo das eras: $+0,925$ em CBO e $+2,317$ em RFC nos 17 repositórios pioneiros entre
E1 e E4, e $+0,266$ em CBO e $+0,984$ em RFC nos 453 repositórios do grupo principal
entre E2 e E4. Em contrapartida, a complexidade ponderada por classe (WMC) apresentou
retração ($-0,893$ no grupo pioneiro e $-0,060$ no grupo principal) e a falta de coesão
(LCOM*) diminuiu ($-0,049$ e $-0,009$), acompanhadas de alta estabilidade na hierarquia
de herança (DIT e NOC). Os achados demonstram que a expansão funcional dos sistemas eleva
a interdependência entre componentes sem provocar degradação na coesão interna nem
inflação na complexidade média das classes.

\textbf{Objetivo (ii), Spring e Boot.} As transições declaradas para dependências do
Spring Boot entre eras adjacentes não exibiram direção estrutural uniforme nas métricas
CK, evidenciando que a adoção do framework opera em interação com o desenho modular de
cada sistema, sem atuar isoladamente como fator determinístico de simplificação ou
complexidade.

\textbf{Objetivo (iii), evidências de IA.} Em E4, entre os 470 repositórios válidos
finais, os 75 sistemas com evidência detectada de assistência por IA apresentaram médias
superiores nas sete medidas estruturais em relação aos 395 sistemas sem evidência
detectada. Contudo, esses projetos possuíam porte médio substancialmente superior
(113.864 contra 7.529 LOC) e já registravam acoplamento mais elevado na era anterior E3
(CBO de 5,858 contra 5,542 nos mesmos grupos), indicando que a presença de artefatos de
agentes associa-se à escala e à intensidade de evolução recente dos sistemas, sem
caracterizar efeito causal direto sobre a arquitetura.

\textbf{Objetivo (iv), curadoria, estrelas e porte.} A exigência de presença contínua de
métricas CK eliminou variações de composição entre eras tanto no grupo pioneiro ($n=17$
em todas as comparações de E1--E4) quanto no grupo principal ($n=453$ em E2 vs.\ E3, E3
vs.\ E4 e E2 vs.\ E4). Demonstrou-se empiricamente que projetos com menos de dez estrelas
validados pelo protocolo (44,5\% da amostra final) apresentam coesão, acoplamento e
tamanho de métodos equivalentes aos de repositórios altamente estrelados, confirmando a
adequação de critérios de engenharia em lugar de métricas de popularidade social.

Como contribuição central, o estudo entrega um protocolo longitudinal auditável e
reprodutível para mineração do ecossistema Spring Boot, aliando curadoria local em quatro
dimensões adaptada às convenções do framework, completude longitudinal de snapshots por
era e análise estrutural sobre conjuntos fixos de projetos. Trabalhos futuros podem
investigar o papel específico de estereótipos arquiteturais do Spring sobre o acoplamento
e aprofundar a rastreabilidade de código assistido por IA em granularidade de método e
pull request, observados os limites descritos na Seção~\ref{sec:ameacas}.
